\section{Conclusion}

% In this study, we conduct the first systematic investigation into the failure modes of LLM-based Multi-Agent Systems (MAS). We analyze over \numDepthTraces execution traces using Grounded Theory, iteratively refining and validating our taxonomy via inter-annotator agreement studies. We identify \numFailureMode fine-grained failure modes, organized into \numFailureArea distinct categories, forming the Multi-Agent System Failure Taxonomy (\taxonomy{}). \taxonomy{} provides a foundational framework for future MAS research. We also develop and validate an automatic evaluation pipeline, LLM Annotator, for scalable failure analysis using \taxonomy{}. This automated annotator serves as a practical tool for developers, enabling systematic diagnosis and evaluation to guide the development of more robust systems.

% We are excited about the potential of MAS, but widespread adoption requires these systems to function reliably. \taxonomy{} contributes towards this goal by providing a framework to understand and mitigate failures. By defining these challenges, we also open concrete problems for the research community to address collaboratively.

In this study, we conduct the first systematic investigation into why MAS fail. This investigation results in the \dataset{}: a comprehensive public resource of over 1600 annotated execution traces from \numMASDomain popular MAS frameworks, which we create to outline MAS failure dynamics and guide future system development. To enable \dataset{}'s systematic annotation and analysis, we first develop the \textbf{M}ulti-\textbf{A}gent \textbf{S}ystem Failure \textbf{T}axonomy (\taxonomy{}). We build \taxonomy{} through a rigorous Grounded Theory-based analysis of an initial 150 traces, validating its definitions with strong inter-annotator agreement and identifying \numFailureMode distinct failure modes across \numFailureArea categories. For scalable annotation of \dataset{} using \taxonomy{}, we then develop an LLM annotator, confirming its high agreement with human experts. Together, \dataset{} and \taxonomy{} provide a foundational framework and empirical grounding for future MAS research.

We are excited about the potential of MAS, but their widespread adoption hinges on achieving greater reliability. Our work, through the public release of \dataset{}, \taxonomy{}, and the LLM annotator, contributes towards this goal. \dataset{} offers a rich empirical basis for understanding current failure dynamics, while \taxonomy{} provides a standardized language and framework to diagnose and mitigate these failures. By systematically identifying and categorizing challenges, we aim to open up concrete research directions and equip the community to develop more robust and effective multi-agent systems.